\documentclass[10pt]{beamer}
\usetheme{metropolis}
\usepackage{graphicx}
\usepackage{listings}
\usepackage{booktabs}
\usepackage{xcolor}
\definecolor{codebg}{HTML}{F5F5F5}
\definecolor{codegreen}{HTML}{2E7D32}
\definecolor{codegray}{HTML}{757575}
\definecolor{codeblue}{HTML}{1565C0}
\lstdefinestyle{rstyle}{language=R,backgroundcolor=\color{codebg},basicstyle=\ttfamily\scriptsize,
  keywordstyle=\color{codeblue}\bfseries,stringstyle=\color{codegreen},
  commentstyle=\color{codegray}\itshape,breaklines=true,frame=single,
  rulecolor=\color{codegray!30},showstringspaces=false,tabsize=2,
  morekeywords={library,ggplot,aes,geom_tile,geom_text,scale_fill_gradient2,
    GGally,ggpairs,ggcorrplot,corrplot,cor,cor.test,chisq.test,
    theme_minimal,labs,lower,upper,diag}}
\lstdefinestyle{pythonstyle}{language=Python,backgroundcolor=\color{codebg},basicstyle=\ttfamily\scriptsize,
  keywordstyle=\color{codeblue}\bfseries,stringstyle=\color{codegreen},
  commentstyle=\color{codegray}\itshape,breaklines=true,frame=single,
  rulecolor=\color{codegray!30},showstringspaces=false,tabsize=4,
  morekeywords={import,from,as,pd,np,plt,sns,heatmap,pairplot,corr,
    pearsonr,spearmanr,kendalltau,chi2_contingency,cramers_v}}
\title{Correlation \& Association Visualization}
\subtitle{Workshop 4 --- Module 6}
\author{Prof.Asoc.\ Endri Raco}
\institute{Data Visualization for Data Scientists \\ UNYT Tirana}
\date{2025/26}
\begin{document}

\begin{frame}
  \titlepage
\end{frame}

\begin{frame}{Module Roadmap}
  \begin{enumerate}
    \item Anscombe's quartet: why scatter before correlate
    \item Three correlation coefficients: Pearson, Spearman, Kendall
    \item Correlation heatmaps: design principles
    \item Pairs plots / scatterplot matrices
    \item Categorical association: Cram\'{e}r's V and mosaic plots
    \item Correlation $\neq$ causation
  \end{enumerate}
  \begin{exampleblock}{Learning Outcome}
    You will always produce a scatter plot before computing a correlation
    coefficient, build publication-quality correlation heatmaps and pairs
    plots, and correctly visualize categorical association.
  \end{exampleblock}
\end{frame}

\section{Anscombe's Quartet}

\begin{frame}{Same $r$, Completely Different Stories}
  \begin{center}
    \includegraphics[width=0.82\textwidth]{figures_w04m06/anscombe}
  \end{center}
\end{frame}

\begin{frame}{The Lesson}
  All four datasets have:
  \begin{itemize}
    \item $r = 0.82$, slope $= 0.50$, intercept $= 3.00$
    \item Same mean of $x$, same mean of $y$, same $R^2$
  \end{itemize}
  \vspace{4pt}
  But only Dataset I is actually linear. II is curved (needs a
  quadratic); III has one outlier driving the fit; IV has one
  leverage point creating the illusion of correlation.
  \begin{alertblock}{Pro-Tip}
    \textbf{Rule}: always scatter plot before computing any correlation.
    A single number ($r$) cannot capture curvature, outliers, clusters,
    or leverage. Anscombe (1973) designed these four datasets specifically
    to make this point.
  \end{alertblock}
\end{frame}

\section{Three Correlation Coefficients}

\begin{frame}{Pearson, Spearman, Kendall}
  \begin{center}
    \includegraphics[width=0.95\textwidth]{figures_w04m06/three_correlations}
  \end{center}
\end{frame}

\begin{frame}{When to Use Which}
  \centering
  \begin{tabular}{lp{4cm}l}
    \toprule
    \textbf{Coefficient} & \textbf{Measures} & \textbf{Use When} \\
    \midrule
    Pearson $r$ & Linear association & Both variables continuous, roughly normal, no outliers \\
    Spearman $\rho$ & Monotonic association (rank-based) & Ordinal data, skewed, or nonlinear but monotonic \\
    Kendall $\tau$ & Concordance of pairs & Small $n$, many ties, or robustness needed \\
    \bottomrule
  \end{tabular}
  \vspace{6pt}

  For the \textbf{U-shaped} relationship (right panel), all three
  return $\approx 0$ --- none captures non-monotonic association.
  Use \textbf{distance correlation} (dcor) or \textbf{mutual information}
  for arbitrary dependence.
\end{frame}

\begin{frame}[fragile]{Computing Correlations}
  \begin{columns}[T]
    \begin{column}{0.48\textwidth}
      \textbf{R}
\begin{lstlisting}[style=rstyle]
# Full matrix
cor(df[, num_cols])           # Pearson
cor(df[, num_cols],
    method = "spearman")      # Spearman

# Single pair with p-value
cor.test(df$x, df$y)
cor.test(df$x, df$y,
    method = "kendall")

# Robust: winsorised
library(WRS2)
wincor(df$x, df$y, tr = 0.1)

# Chi-squared for categorical
chisq.test(table(df$cat1,
                  df$cat2))
\end{lstlisting}
    \end{column}
    \begin{column}{0.48\textwidth}
      \textbf{Python}
\begin{lstlisting}[style=pythonstyle]
# Full matrix
df[num_cols].corr()          # Pearson
df[num_cols].corr(
    method="spearman")       # Spearman

# Single pair with p-value
from scipy.stats import (
    pearsonr, spearmanr,
    kendalltau)
r, p = pearsonr(x, y)
rho, p = spearmanr(x, y)

# Chi-squared for categorical
from scipy.stats import (
    chi2_contingency)
ct = pd.crosstab(df.cat1, df.cat2)
chi2, p, dof, exp = (
    chi2_contingency(ct))
\end{lstlisting}
    \end{column}
  \end{columns}
\end{frame}

\section{Correlation Heatmaps}

\begin{frame}{Five Design Principles}
  \begin{center}
    \includegraphics[width=0.92\textwidth]{figures_w04m06/corrplot_principles}
  \end{center}
\end{frame}

\begin{frame}{Annotated Heatmap}
  \begin{center}
    \includegraphics[width=0.62\textwidth]{figures_w04m06/heatmap}
  \end{center}
\end{frame}

\begin{frame}[fragile]{Heatmap in Code}
  \begin{columns}[T]
    \begin{column}{0.48\textwidth}
      \textbf{R (ggcorrplot or corrplot)}
\begin{lstlisting}[style=rstyle]
library(ggcorrplot)

corr <- cor(df[, num_cols],
    use = "pairwise.complete")

ggcorrplot(corr,
  type = "lower",        # triangle
  lab = TRUE,            # annotate
  lab_size = 3,
  colors = c("#E53935",  # neg
    "white",             # zero
    "#1565C0"),          # pos
  hc.order = TRUE,       # cluster
  outline.color = "white")

# Alternative: corrplot package
library(corrplot)
corrplot(corr, method = "color",
  type = "lower", tl.col = "black",
  addCoef.col = "black",
  order = "hclust")
\end{lstlisting}
    \end{column}
    \begin{column}{0.48\textwidth}
      \textbf{Python (seaborn)}
\begin{lstlisting}[style=pythonstyle]
corr = df[num_cols].corr()

# Lower triangle mask
mask = np.triu(np.ones_like(
    corr, dtype=bool))

sns.heatmap(
    corr,
    mask=mask,             # triangle
    annot=True,            # annotate
    fmt=".2f",
    cmap="RdBu_r",        # diverging
    vmin=-1, vmax=1,       # fixed scale
    center=0,
    square=True,
    linewidths=0.5,
    cbar_kws={"shrink": 0.8,
        "label": "Pearson r"})
\end{lstlisting}
    \end{column}
  \end{columns}
\end{frame}

\section{Pairs Plots}

\begin{frame}{The Scatterplot Matrix}
  \begin{center}
    \includegraphics[width=0.72\textwidth]{figures_w04m06/pairs}
  \end{center}
\end{frame}

\begin{frame}[fragile]{Pairs Plot in Code}
  \begin{columns}[T]
    \begin{column}{0.48\textwidth}
      \textbf{R (GGally)}
\begin{lstlisting}[style=rstyle]
library(GGally)

ggpairs(df,
  columns = c("Rating",
    "log_reviews", "Size"),
  mapping = aes(color = Type,
    alpha = 0.3),
  lower = list(continuous =
    wrap("points", size = 0.5)),
  diag = list(continuous =
    wrap("densityDiag")),
  upper = list(continuous =
    wrap("cor", size = 3))) +
  theme_minimal(base_size = 8)
\end{lstlisting}
    \end{column}
    \begin{column}{0.48\textwidth}
      \textbf{Python (seaborn)}
\begin{lstlisting}[style=pythonstyle]
sns.pairplot(
    df[["Rating", "log_reviews",
        "Size", "Type"]],
    hue="Type",
    palette={"Free": "#1565C0",
             "Paid": "#E53935"},
    diag_kind="kde",
    plot_kws={"s": 10,
              "alpha": 0.3},
    height=2)
\end{lstlisting}
    \end{column}
  \end{columns}
  \begin{alertblock}{Pro-Tip}
    Pairs plots scale poorly beyond $\sim$8 variables (64 panels).
    For high-dimensional data, use a \textbf{correlation heatmap}
    first to identify the strongest pairs, then drill down with
    individual scatter plots or pairs on a subset.
  \end{alertblock}
\end{frame}

\section{Categorical Association}

\begin{frame}{Visualizing Categorical $\times$ Categorical}
  \begin{center}
    \includegraphics[width=0.92\textwidth]{figures_w04m06/categorical}
  \end{center}
  \begin{exampleblock}{Data Insight}
    For two categorical variables, the visual equivalent of a chi-squared
    test is a \textbf{100\% stacked bar} (proportions) or a
    \textbf{balloon plot} (counts, size $\propto$ count). If all bars
    look identical, the variables are independent. Differences in bar
    composition $\Rightarrow$ association. Quantify with
    \textbf{Cram\'{e}r's V} $\in [0, 1]$.
  \end{exampleblock}
\end{frame}

\section{Correlation $\neq$ Causation}

\begin{frame}{The Most Important Slide in This Module}
  \begin{center}
    \includegraphics[width=0.82\textwidth]{figures_w04m06/not_causation}
  \end{center}
  \begin{alertblock}{Pro-Tip}
    Three sources of spurious correlation: \textbf{confounders} (a third
    variable causes both), \textbf{reverse causation} ($Y$ causes $X$,
    not $X$ causes $Y$), and \textbf{coincidence} (random co-trend in
    finite samples). Correlation establishes association; only controlled
    experiments or careful causal inference (DAGs, instrumental variables)
    establish causation.
  \end{alertblock}
\end{frame}

\section{Synthesis}

\begin{frame}{Key Takeaways}
  \begin{enumerate}
    \item \textbf{Always scatter first}. Anscombe's quartet proves that
          $r$ alone is meaningless without a visual check.
    \item \textbf{Pearson} = linear; \textbf{Spearman} = monotonic;
          \textbf{Kendall} = concordance. None captures non-monotonic
          dependence --- use distance correlation for that.
    \item Heatmap design: \textbf{diverging palette, annotate cells,
          cluster rows/columns, lower triangle only, fixed $[-1,1]$
          colour scale}.
    \item \textbf{Pairs plots} combine scatter + density + $r$ in one
          figure. Limit to $\leq 8$ variables.
    \item For categorical association: \textbf{100\% stacked bar} or
          \textbf{balloon plot} + Cram\'{e}r's V.
    \item \textbf{Correlation $\neq$ causation}. Always consider
          confounders, reverse causation, and coincidence.
  \end{enumerate}
\end{frame}

\begin{frame}{Essential Readings}
  \begin{itemize}
    \item Anscombe, F.~J. (1973). ``Graphs in Statistical Analysis.''
          \emph{The American Statistician}, 27(1), 17--21.
    \item Wilke, C.~O. (2019). \emph{Fundamentals of Data Visualization},
          Chapter 12 (Visualizing Associations).
    \item Friendly, M. (2002). ``Corrgrams.'' \emph{The American
          Statistician}, 56(4), 316--324.
    \item Pearl, J. (2009). \emph{Causality}, 2nd ed. Cambridge
          University Press. (Correlation vs causation.)
  \end{itemize}
  \vspace{6pt}
  \textbf{Next Module}: EDA Case Study --- Google Play Store (W04-M07)
\end{frame}

\end{document}
